% SOURCE_STATEMENT: "For each fixed rank, are there only finitely many fusion categories, and likewise only finitely many ribbon categories, up to tensor equivalence?"
% Preparation and final self-audit: CONTRIBUTING.md prompts C and D; source/audit artifacts retained in the local physics_statements_2026_10_11 batch.
% !TeX program = lualatex
\documentclass[11pt]{article}
\usepackage[margin=25mm]{geometry}
\usepackage{fontspec}
\setmainfont{Latin Modern Roman}
\usepackage{amsmath,amssymb,mathtools}
\usepackage{unicode-math}
\setmathfont{Latin Modern Math}
\usepackage{xurl,hyperref}
\hypersetup{colorlinks=true,urlcolor=blue}
\setcounter{secnumdepth}{0}
\setlength{\parindent}{0pt}
\setlength{\parskip}{5pt}
\setlength{\emergencystretch}{3em}
\newcommand{\sref}[2]{\hyperlink{source:#1:#2}{[S#2]}}
\newcommand{\cataloguescope}[1]{\paragraph{Recorded target.}#1}
\newcommand{\R}{\mathbb R}
\newcommand{\C}{\mathbb C}
\newcommand{\N}{\mathbb N}
\newcommand{\Z}{\mathbb Z}
\newcommand{\Q}{\mathbb Q}
\newcommand{\F}{\mathbb F}
\newcommand{\eps}{\varepsilon}
\begin{document}
\section{Wang–Ostrik Finiteness Conjecture for Fusion Categories}\label{statement}

\subsection{Definitions and mathematical statement}

\paragraph{Shared notation.}$\mathbb N=\{1,2,\ldots\}$; the standard number fields have their usual meanings. All additional conventions are stated locally.


\paragraph{Definitions and conventions.} Over $\mathbb C$, a fusion category is a finite semisimple rigid linear tensor category with simple tensor unit and finite-dimensional morphism spaces. Its rank is the number of isomorphism classes of simple objects. A ribbon category in this question is a fusion category equipped with the compatible braiding, duality, and twist structure. Tensor equivalence identifies categories by a linear monoidal equivalence, with the relevant additional structures respected for the ribbon version.

\paragraph{Source scope and qualifications.} The fixed-rank parameter ranges over positive integers. The report distinguishes the characteristic-zero setting from positive characteristic, and its rank-two and ribbon-rank-three results are special cases rather than proof of the full conjecture.

\paragraph{Catalogue mathematical target.}
\paragraph{Catalogue wording retained for traceability.}
For each fixed rank, are there only finitely many fusion categories, and likewise only finitely many ribbon categories, up to tensor equivalence?

\paragraph{Saved literature context.} The supplied catalogue's assessment dated 2026-08-21 records: The Wang--Ostrik finiteness conjecture for all fusion categories of a fixed rank remains open. Rank-finiteness is known for important subclasses such as modular and braided fusion categories. The assessment is retained as dated provenance; it is not a new comprehensive literature or open-status certification.
\subsection{Short English statement}
For each prescribed number of simple objects, decide whether only finitely many fusion and ribbon categories can occur.
\subsection{Sources}
\begin{itemize}
\item[\hypertarget{source:statement:1}{S1}] ULAM.AI and the UnsolvedMath Contributors, UnsolvedMath, catalogue entry. Catalogue curation under CC BY 4.0; definitions and English exposition added locally.
\url{https://huggingface.co/datasets/ulamai/UnsolvedMath}.
\item[\hypertarget{source:statement:2}{S2}] Arbeitsgemeinschaft: Conformal Field Theory (2007), Oberwolfach Reports, DOI 10.4171/owr/2007/17, pp. 945--994; official report source and mathematical context.
\url{https://doi.org/10.4171/owr/2007/17}.
\item[\hypertarget{source:statement:3}{S3}] Official EMS Press PDF of the same Oberwolfach report; PDF page 23.
\url{https://ems.press/content/serial-article-files/46103}.
\item[\hypertarget{source:statement:4}{S4}] rank-finiteness for modular categories; reference in the saved dated catalogue literature review.
\url{https://arxiv.org/abs/1310.7050}.
\end{itemize}
\label{end:statement}
\end{document}
